\maketitle
\vspace{-1.5em}
{\small\noindent $^{*}$ These authors contributed equally to this work.\\
$^{**}$ Corresponding author: \href{mailto:pierre-henri.dubuis@agroscope.admin.ch}{pierre-henri.dubuis@agroscope.admin.ch}\par}

\begin{abstract}
Agroscope conducted a long-term survey to assess the incidence of esca disease in Western Switzerland between 2001 and 2022 and to characterize its temporal dynamics. This study has four main objectives: (1) to provide a rare and particularly valuable long-term dataset documenting the incidence of esca over more than two decades; (2) to identify potentially important climatic drivers using a recent model-independent variable importance algorithm that quantifies the strength of their association with esca incidence; (3) to develop predictive models of esca incidence using climatic covariates derived from microclimatic weather data; and (4) to investigate the relationship between water balance and regionally aggregated esca incidence. The climatic variables identified by the variable importance algorithm as being most strongly associated with esca incidence at the plot level were primarily related to water balance.

The predictive analysis showed that model performance at the plot level was limited. Consistent with previous studies, analysis of the variables contributing to model predictions did not allow us to identify a single factor, or set of factors, that could reliably predict esca incidence.

At the regional scale, however, water balance and seasonal vapor pressure deficit (VPD) were strongly associated with esca incidence across years, with drier regions generally showing lower disease incidence.

Finally, addressing the initial question that motivated this long-term survey, we observed a slight decreasing trend in esca incidence over the study period. This trend coincided with increasingly dry conditions in Western Switzerland and may therefore be partly associated with changes in regional water balance in the context of climate change.

All the datasets and analyses presented here are made available and reproducible.
\end{abstract}

\noindent\textbf{Keywords:} esca disease, predictive modelling, FORD algorithm
